\section{From bounded blocks to AP without BAP}\label{sec:reduction}

The following sections will construct the blocks of \Cref{thm:blocks}.
Here we assume that theorem, prove the renorming and wedge results, and
derive the nonlinear consequence from Kalton's theorem. A bounded space with separation at least one has a
coefficient description by $\ell_1$, with an approximation bound
controlled by its diameter.
Joining the blocks preserves each of their obstructions, while retractions
onto finitely many blocks provide the approximations required for AP.

\subsection{The coefficient estimate}

The coefficient identification below appears in
\cite[Theorem~4.14]{AACD20}; the coefficient bounds are stated in
\cite[Lemma~3.4]{CuthKalenda26}. We include the proof and derive the
stated approximation bound from coordinate truncations.

\begin{proposition}\label{prop:coefficients}
Let $(N,d,o)$ be a countable pointed metric space such that
$d(x,y)\ge1$ for distinct points and $\operatorname{diam}N\le B<\infty$.
The map
\[
 J:\ell_1(N\setminus\{o\})\longrightarrow\F(N),
 \qquad Jc=\sum_{x\ne o}c_x\delta_x,
\]
is a linear isomorphism onto $\F(N)$ and satisfies
\begin{equation}\label{eq:coefficient-bounds}
 \tfrac12\norm{c}_1\le\norm{Jc}\le B\norm{c}_1.
\end{equation}
If $N\ne\{o\}$, then $\F(N)$ has $2B$-BAP. If $N=\{o\}$, both
spaces are zero.
\end{proposition}

\begin{proof}
First let $c$ be finitely supported. Since $\norm{\delta_x}=d(x,o)\le B$,
the triangle inequality gives the upper estimate in
\eqref{eq:coefficient-bounds}. For the lower estimate, define
$f_c:N\to\R$ to be $\operatorname{sign}(c_x)/2$ at each nonbasepoint
in the support of $c$, and zero elsewhere. This function vanishes at
$o$. Its values differ by at most one, so the separation assumption
makes it $1$-Lipschitz. Testing the evaluation combination against $f_c$
gives
\[
 \norm{Jc}\ge (Jc)(f_c)=\tfrac12\sum_{x\ne o}|c_x|.
\]

The upper estimate makes the series defining $Jc$ converge in $\F(N)$
for every $c\in\ell_1(N\setminus\{o\})$. Both estimates pass to this
extension by approximation with finitely supported families. In particular,
$J$ is injective and its inverse on its range has norm at most two.
The range is closed: if $Jc_m$ converges, then
\[
 \norm{c_m-c_j}_1\le2\norm{Jc_m-Jc_j}
\]
makes $(c_m)$ Cauchy, and continuity of $J$ identifies the limit of
$(Jc_m)$ as the image of the limit of $(c_m)$. The range also contains
every evaluation $\delta_x$, whose span is dense in $\F(N)$. It is
therefore all of $\F(N)$.

Suppose now that $N\ne\{o\}$; then $B\ge1$. Choose increasing finite
sets $E_m$ exhausting $N\setminus\{o\}$, and let $q_m$ be the associated
coordinate projections on $\ell_1(N\setminus\{o\})$. The operators
\[
 S_m=Jq_mJ^{-1}:\F(N)\longrightarrow\F(N)
\]
have finite rank and norm at most $B\cdot1\cdot2=2B$. Since
$q_mc\to c$ in $\ell_1$, they satisfy $S_mu\to u$ for every
$u\in\F(N)$. This convergence is uniform on compact sets. Indeed, for
a nonempty compact $C\subseteq\F(N)$ and $\varepsilon>0$, choose a
finite $\varepsilon/[3(2B+1)]$-net in $C$. For all sufficiently large
$m$, the error $\norm{S_mu-u}$ is less than $\varepsilon/3$ at every
point of the net. Since $\norm{S_m-\Id}\le2B+1$, the error at any
point of $C$ is then less than $2\varepsilon/3$, hence less than
$\varepsilon$. The empty compact set imposes no condition. If $N$ is
finite, the sequence of coordinate projections is eventually the identity.
This proves $2B$-BAP. The assertion for $N=\{o\}$ is immediate.
\end{proof}

\begin{proof}[Proof of \Cref{cor:renorming}]
Fix $p\ge1$. By \Cref{thm:blocks}, the set $M_p\setminus\{o\}$ is
countably infinite, so an enumeration identifies its coefficient space
with real $\ell_1$. Apply \Cref{prop:coefficients} with $B=B_p$, denote
its coefficient map by $J_p$, and define
$\norm{c}_p=\norm{J_pc}_{\F(M_p)}$. Equation
\eqref{eq:coefficient-bounds} gives the stated norm inequalities and
shows that this is an equivalent norm. For this norm, $J_p$ is a
surjective isometry. Conjugation by $J_p$ preserves operator norms,
finite rank, and approximation errors on corresponding compact sets.
It therefore transfers $2B_p$-BAP to the new norm, with
$2B_p=600p^2+300p+10$. If the new norm had $p$-BAP, the same conjugation
would give $p$-BAP on $\F(M_p)$. Applying it with accuracy $1/2$ to
the finite set $\{\delta_a:a\in A_p\}$ would contradict
\Cref{thm:blocks}. For $p=1$, this is precisely failure of MAP.
\end{proof}

\subsection{The metric wedge and its retractions}

Take disjoint copies of the spaces $M_p$ and identify only their
basepoints. Thus, as a set,
\[
 M=\{o\}\cup\bigsqcup_{p\ge1}(M_p\setminus\{o\}).
\]
Keep the distance $d_p$ within each block, and for nonbasepoint points
in distinct blocks set
\begin{equation}\label{eq:wedge-metric}
 d(x,y)=d_p(x,o)+d_q(y,o)
 \qquad(x\in M_p\setminus\{o\},\ y\in M_q\setminus\{o\},\ p\ne q).
\end{equation}
These rules define a metric. Symmetry and positivity are immediate.
For the triangle inequality, if the endpoints lie in one block, a
detour through a point $z$ in another block has length
$d_p(x,o)+2d_q(z,o)+d_p(y,o)\ge d_p(x,y)$. If the endpoints lie in
different blocks and the third point lies in one of those blocks, that
block's triangle inequality gives the required bound. A third point in
a third block adds twice its distance to $o$, and a third point equal
to $o$ gives equality. The remaining case is the triangle inequality
within a single block.

The space $M$ is countable. Within a block, distinct points remain at
distance at least one; in distinct blocks, nonbasepoint points are at
distance at least two. It is also unbounded. To see this using only
\Cref{thm:blocks}, suppose that $M$ had finite diameter $C_0$. Then
$C_0\ge1$, and every $M_p$ would have diameter at most $C_0$.
\Cref{prop:coefficients} would give $2C_0$-BAP on every $\F(M_p)$.
Choose an integer $p\ge\max\{1,2C_0\}$. Applying that property to
$\{\delta_a:a\in A_p\}$ with accuracy $1/2$ would produce a finite-rank
operator of norm at most $2C_0\le p$ whose error is less than $1/2$
at every anchor, contrary to \Cref{thm:blocks}.

Finally, the closed ball about $o$ of radius $B_1$ contains the
countably infinite $1$-separated set $M_1$. If this ball were compact,
its cover by open balls of radius $1/3$ would admit a finite subcover.
Each such ball contains at most one point of $M_1$, which is impossible.
Thus $M$ is nonproper.

For each $p$, let $i_p:M_p\to M$ be inclusion and let $r_p:M\to M_p$
fix $M_p$ and send all other blocks to $o$. Both are pointed
$1$-Lipschitz maps. Inclusion preserves distances. For the retraction,
the only case that needs an inequality has $x\in M_p$ and
$y\in M_q\setminus\{o\}$ for $q\ne p$; then
\[
 d_p(r_p(x),r_p(y))=d_p(x,o)\le d(x,y).
\]
If both points are outside $M_p$, their images coincide, and if both
are in $M_p$, their distance is preserved.

We use the elementary linearization property of free spaces; its role
in Banach-space geometry is developed in \cite[\S2]{GodefroyKalton03}.
If $u:N\to N'$ is pointed and $1$-Lipschitz, define on finite evaluation
combinations
\[
 \widehat u\left(\sum_x c_x\delta_x\right)=\sum_x c_x\delta_{u(x)}.
\]
Pullback of pointed scalar tests gives
\begin{align*}
 \left\|\sum_x c_x\delta_{u(x)}\right\|
 &=\sup_{\substack{f\in\Lip_0(N')\\\Lipnorm{f}\le1}}
       \left|\sum_x c_x(f\circ u)(x)\right|\\
 &\le\left\|\sum_x c_x\delta_x\right\|.
\end{align*}
This inequality makes the definition independent of the representation
of the evaluation combination, and it gives a unique continuous linear
extension $\widehat u:\F(N)\to\F(N')$ of norm at most one.
Applying it to $i_p$ and $r_p$ gives contractions
\[
 I_p:\F(M_p)\longrightarrow\F(M),\qquad
 P_p:\F(M)\longrightarrow\F(M_p).
\]
Their composite $P_pI_p$ fixes every evaluation, so it is the identity
on $\F(M_p)$. In particular, for every $v\in\F(M_p)$,
\[
 \norm{v}=\norm{P_pI_pv}\le\norm{I_pv}\le\norm{v};
\]
thus $I_p$ is an isometry.

\subsection{The approximation properties of the wedge}

Suppose that $\F(M)$ had $\lambda$-BAP for some finite $\lambda\ge1$.
Choose an integer $p\ge\lambda$ and apply that property with accuracy
$1/2$ to the finite, hence compact, set
$\{I_p\delta_a:a\in A_p\}$. Let $U$ be the resulting operator.
The compression $P_pUI_p$ has finite rank and norm at most
$\lambda\le p$, while for every $a\in A_p$,
\[
 \norm{P_pUI_p\delta_a-\delta_a}
 \le\norm{UI_p\delta_a-I_p\delta_a}<\tfrac12.
\]
This contradicts \Cref{thm:blocks}. All blocks and their anchor sets
were fixed before $\lambda$ and $U$ were chosen.

It remains to prove AP. Kalton's general theorem
\cite[Proposition~4.4]{Kalton04} gives AP for free spaces over uniformly
discrete metric spaces; the wedge structure provides a direct proof here.
For $N\ge1$, define
$r_{\le N}:M\to M$ to retain the first $N$ blocks and send every other
block to $o$. It is pointed and $1$-Lipschitz: distances in the retained
part are unchanged, distances in the collapsed part become zero, and
\eqref{eq:wedge-metric} bounds the image distance for a retained point
and a collapsed point. Its linearization is a contractive projection
\begin{equation}\label{eq:finite-block-projection}
 R_N=\sum_{p=1}^N I_pP_p.
\end{equation}
The equality follows on each evaluation $\delta_x$: the sum retains it
when $x$ is in one of the first $N$ blocks and sends it to zero otherwise.
The same check gives $R_N^2=R_N$. Continuity gives both identities on
$\F(M)$.

Every finite linear combination of evaluations is fixed by $R_N$ for
all sufficiently large $N$. These combinations are dense, and
$\norm{\Id-R_N}\le2$. Thus $R_Nu\to u$ for every $u\in\F(M)$:
given a finite evaluation combination $v$ fixed by $R_N$, one has
$\norm{R_Nu-u}\le2\norm{u-v}$. The finite-net argument in the proof of
\Cref{prop:coefficients}, applied now with $\norm{\Id-R_N}\le2$, makes
this convergence uniform on every norm-compact set.

Now fix a nonempty norm-compact $C\subseteq\F(M)$ and
$\varepsilon>0$. First choose $N\ge1$ such that
\[
 \sup_{u\in C}\norm{R_Nu-u}<\varepsilon/2.
\]
For each $1\le p\le N$, the image $P_p(C)$ is compact. By
\Cref{prop:coefficients}, choose a finite-rank operator
$T_p:\F(M_p)\to\F(M_p)$ with $\norm{T_p}\le2B_p$ and
\[
 \sup_{u\in C}\norm{T_pP_pu-P_pu}<\frac{\varepsilon}{2N}.
\]
The operator
\[
 T=\sum_{p=1}^N I_pT_pP_p:\F(M)\longrightarrow\F(M)
\]
has finite rank and is bounded, with
$\norm{T}\le\sum_{p=1}^N2B_p<\infty$. By
\eqref{eq:finite-block-projection}, for every $u\in C$,
\[
 \norm{Tu-u}
 \le\norm{R_Nu-u}
       +\sum_{p=1}^N\norm{T_pP_pu-P_pu}<\varepsilon.
\]
We first chose $N$ so that the contractive projection $R_N$ approximated
the whole compact set, and then chose the finitely many block operators.
Their combined norm may depend on $C$ and
$\varepsilon$, as AP permits. This proves AP, and the metric properties
and the preceding obstruction prove \Cref{thm:main} under
\Cref{thm:blocks}.

\subsection{Failure of nonlinear approximability}\label{sec:nonlinear}

We verify the nonlinear consequence stated in the introduction using
Kalton's characterization by nets \cite[Theorem~4.4 and p.~1260]{Kalton12}.
Let $X=\F(M)$ for the
wedge just constructed. It is separable because $M$ is countable, and
$\{\delta_x:x\in M\}$ is $1$-separated in $X$. Extend this set to a maximal
$1$-separated set $N_X\subseteq X$. Every point of $X$ is then within
distance one of $N_X$, so $N_X$ is a net: a separated set within a fixed
distance of every point of $X$. If $X$ were approximable, Kalton's theorem
would give one Lipschitz bound $\Lambda\ge1$ for maps $g:N_X\to N_X$
with relatively compact range that approximate any prescribed finite subset
as closely as desired. Separation makes each such range finite and turns
approximation within distance less than one into exact fixing. For a finite
$F\subseteq M$ containing $o$, choose $g$ fixing every $\delta_x$ with
$x\in F$. The restriction $x\mapsto g(\delta_x)$ is then a pointed
$\Lambda$-Lipschitz map from $M$ into $X$. The Banach-valued linearization
property of free spaces \cite[Lemma~3.2]{Kalton04} gives a finite-rank
operator $T_F:X\to X$ of norm at most $\Lambda$ satisfying
$T_F\delta_x=g(\delta_x)$ for all $x\in M$. To approximate a compact set,
approximate a finite net in it by finite evaluation combinations and put
all their supports in $F$. The operator $T_F$ fixes those combinations,
and $\norm{T_F-\Id}\le\Lambda+1$ controls the remaining errors. Thus these
operators give BAP, contrary to the preceding obstruction. Hence $X$ is not
approximable.

It remains to construct the blocks of \Cref{thm:blocks}.
